\documentclass[12pt, a4paper]{article}

\usepackage[utf8]{inputenc}
\usepackage[T1]{fontenc}
\usepackage{geometry}
\usepackage{hyperref}
\usepackage{amsmath}
\usepackage{booktabs}
\usepackage{amssymb}
\usepackage[backend=biber, style=authoryear]{biblatex}

\geometry{margin=2.5cm}
\addbibresource{references.bib}

\title{Image Fragment Reconstruction via Self-Supervised Learning\\
\large Task Summary and Algorithm Design}
\author{}
\date{\today}

\begin{document}

\maketitle
\tableofcontents
\newpage

% ─────────────────────────────────────────────
\section{Task Recapitulation}
% ─────────────────────────────────────────────

\subsection{Core Objective}

The task is to develop a self-supervised machine learning model that groups a large
collection of mixed image fragments back to their original source images.
Concretely, the pipeline is as follows:

\begin{enumerate}
    \item A data generator yields a sample of \textbf{10 images} (64$\times$64$\times$3 pixels each),
          with augmentation optionally applied before delivery.
    \item Each image is cut into a \textbf{4$\times$4 grid} of non-overlapping tiles,
          producing 16 fragments of 16$\times$16$\times$3 pixels per image.
    \item All 160 fragments (10 images $\times$ 16 tiles) are \textbf{shuffled into a single
          unordered collection} — no label, no position information is retained.
    \item The model must learn to \textbf{group these 160 fragments back into 10 clusters},
          each corresponding to one source image.
\end{enumerate}

\subsection{Points That Required Clarification}

Several aspects of the task were not immediately obvious and deserve explicit statement:

\paragraph{Self-supervised does not mean label-free throughout.}
The ground-truth grouping (which fragment came from which image) is known to the
experimenter at all times — it is used to construct training pairs and to evaluate the
model. What the model never receives is an explicit label. The supervision signal
is derived from the data generation process itself, not from human annotation.

\paragraph{The goal is grouping, not spatial reconstruction.}
The model is not required to determine where each fragment was positioned within the
original image (top-left, bottom-right, etc.). It only needs to identify which fragments
originated from the same image. Think of sorting puzzle pieces back into their correct
boxes without assembling the puzzles.

\paragraph{The generator yields augmented images, not augmented fragments.}
Augmentation is applied to the full 64$\times$64 image \emph{before} it is cut into
fragments. Consequently, all 16 fragments derived from one generator call share
a consistent (augmented) visual appearance. The augmentations applied are:
rotation (up to 20°), horizontal and vertical shift (up to 10\%), and channel shift
(up to 0.2). These are defined in \texttt{data.py} via \texttt{init\_augmentor()}.

\paragraph{Augmentation as an implicit constraint.}
Because augmentation randomly perturbs colours and positions, any model that relies on
exact pixel values or absolute spatial coordinates will fail. The model is therefore
forced to learn higher-level features — texture, pattern, and visual style — that
remain stable across an image regardless of minor transformations.

\paragraph{The number of clusters is known.}
Since the experimenter controls how many images are fed into each sample (always 10),
the number of clusters $k$ is known in advance. Algorithms that require $k$ as input
(e.g.\ $k$-means) can therefore be used directly. In a real-world deployment where
$k$ is unknown, discovery-based methods (e.g.\ DBSCAN) would be required.

\paragraph{Balanced clusters.}
Because each image is always cut into exactly 16 fragments, the true clusters are
perfectly balanced. This can be exploited by constraining the clustering step to
produce balanced clusters of 16, which reduces the search space and can improve
accuracy.

\paragraph{Infinite effective training data.}
Since augmentations are sampled randomly on every generator call, no two batches are
identical. The training set is effectively infinite, and the experimenter controls
training duration by choosing how many generator calls to make.

% ─────────────────────────────────────────────
\section{Related Problems}
% ─────────────────────────────────────────────

The task borrows from three well-established areas of machine learning research.

\subsection{Jigsaw Puzzle Solving}

\textcite{noroozi2016} proposed using jigsaw puzzle reassembly as a self-supervised
pretext task: image patches are shuffled and a network is trained to predict their
correct spatial arrangement. This is closely related to the present task in that both
work with fragments of images without labels. The key difference is that jigsaw methods
require recovering the \emph{spatial position} of each fragment, whereas here only
\emph{source identity} (which image did this fragment come from) must be determined.

\subsection{Instance Discrimination and Contrastive Learning}

\textcite{simclr2020} introduced SimCLR, a framework in which two augmented views of the
same image are treated as a positive pair and all other images in the batch as negatives.
The encoder is trained so that positive pairs have similar embeddings and negative pairs
have dissimilar embeddings. This is essentially the training objective used in the present
task, with fragments from the same source image playing the role of positive pairs.

\textcite{moco2020} proposed MoCo, which extends contrastive learning with a momentum
encoder and a memory bank to increase the number of negatives available per update.
MoCo is most beneficial when batch sizes are small. In the present task the number of
negatives per batch is already manageable (at most 144 fragments from other images within
a single 160-fragment sample), making the additional complexity of MoCo unnecessary.

\subsection{Deep Clustering}

Standard unsupervised image clustering methods learn embeddings and cluster assignments
jointly. The present task is distinct in that the unit to be clustered is a
\emph{fragment}, not a whole image, and the cluster identity is ephemeral — it changes
with every generator call and carries no semantic meaning beyond the current sample.

\subsection{Summary of Similarities and Differences}

\begin{center}
\begin{tabular}{lccc}
\toprule
 & \textbf{Jigsaw} & \textbf{SimCLR} & \textbf{This task} \\
\midrule
Works with fragments & \checkmark & & \checkmark \\
Recovers spatial position & \checkmark & & \\
Contrastive objective & & \checkmark & \checkmark \\
Consistent class identity across batches & & \checkmark & \\
Clustering at inference & & & \checkmark \\
\bottomrule
\end{tabular}
\end{center}

% ─────────────────────────────────────────────
\section{Algorithm Outline and Decision Points}
% ─────────────────────────────────────────────

\subsection{Overview}

The proposed pipeline has three stages:

\begin{enumerate}
    \item \textbf{Data preparation} — fragment images and construct training pairs.
    \item \textbf{Training} — learn an encoder that maps fragments to embeddings using
          a contrastive loss.
    \item \textbf{Evaluation} — cluster the embeddings with $k$-means and measure
          grouping accuracy.
\end{enumerate}

\subsection{Stage 1: Data Preparation}

Call the generator with \texttt{batch\_size=10} and \texttt{augmentation=True}. Cut each
of the 10 returned images into a 4$\times$4 grid, yielding 160 fragments of
16$\times$16$\times$3. Retain the source-image index for each fragment — this is used to
construct training pairs and to evaluate, but is never passed to the model.

\subsection{Stage 2: Encoder Architecture}

\paragraph{Decision point: architecture.}
A small \textbf{CNN encoder} is the natural choice for 16$\times$16 image patches.
CNNs extract local spatial features (edges, textures, patterns) efficiently at this
scale. Vision Transformers are designed for larger inputs and would be overkill here.
An autoencoder could also produce embeddings via its bottleneck, but its reconstruction
objective does not directly optimise for inter-group separation.

\paragraph{Decision point: embedding dimension.}
A dimension of 128--256 is standard for contrastive methods and provides enough capacity
without making clustering unnecessarily difficult.

\paragraph{Decision point: projection head.}
Following \textcite{simclr2020}, a two-layer MLP projection head is added on top of the
encoder during training and discarded at evaluation time. This is empirically shown to
improve the quality of the encoder representations.

\subsection{Stage 3: Contrastive Loss}

The NT-Xent loss from SimCLR \parencite{simclr2020} is used. For each fragment $i$, its
positive set consists of the other 15 fragments from the same source image. All remaining
fragments in the batch serve as negatives. The loss encourages:

\[
\text{small intra-group distance} \quad \text{and} \quad \text{large inter-group distance}
\]

\paragraph{Decision point: distance metric.}
\textbf{Cosine similarity} is used as the distance measure, consistent with SimCLR.
Cosine similarity measures the angle between embedding vectors, making it invariant to
the magnitude of the embeddings.

\paragraph{Decision point: training duration.}
Since each generator call produces a unique augmented sample, training duration is
controlled by the number of generator calls (iterations). This is a hyperparameter to
tune; more iterations generally improve representations up to a point of diminishing
returns.

\subsection{Stage 4: Clustering}

\paragraph{Decision point: clustering algorithm.}
\textbf{$k$-means} with $k=10$ is used, since the number of source images per sample is
known. The key constraint is that the clustering algorithm must use the \textbf{same
distance metric as the contrastive loss}. Because cosine similarity was used during
training, embeddings should be $\ell_2$-normalised before clustering — on the unit
hypersphere, Euclidean distance and cosine distance are equivalent, so standard
$k$-means applies correctly.

If Euclidean distance were used during training instead, $\ell_2$-normalisation would
not be appropriate and standard $k$-means could be applied directly to unnormalised
embeddings. Mismatching the training metric and the clustering metric would result in
clusters that do not align with what the model actually learned.

\paragraph{Decision point: balanced clusters.}
Since the true cluster size is always exactly 16, the clustering can be constrained to
produce balanced clusters. This reduces ambiguity and can improve accuracy.

\subsection{Stage 5: Evaluation Metrics}

\paragraph{Decision point: metric choice.}
Simple \textbf{cluster purity} measures the fraction of the dominant class in each
cluster. It is intuitive but ignores how errors are distributed. For example, a cluster
containing 60\% fragments from image A and 40\% from image B scores identically to a
cluster containing 60\% from image A, 20\% from image B, and 20\% from image C —
despite the latter being more disordered.

\textbf{Adjusted Rand Index (ARI)} and \textbf{Normalised Mutual Information (NMI)}
capture the full distribution of fragments across all clusters simultaneously and are
therefore the preferred metrics. Both are standard in clustering evaluation and correct
for chance agreement.

\subsection{Summary of Decision Points}

\begin{center}
\begin{tabular}{lll}
\toprule
\textbf{Decision} & \textbf{Choice} & \textbf{Alternatives} \\
\midrule
Encoder architecture & Small CNN & Autoencoder, ViT \\
Embedding dimension & 128--256 & Smaller / larger \\
Training objective & NT-Xent (SimCLR) & Triplet loss, BYOL \\
Distance metric & Cosine similarity & Euclidean, dot product \\
Clustering algorithm & $k$-means ($k=10$) & DBSCAN, spectral \\
Cluster size constraint & Balanced (16 per cluster) & Unconstrained \\
Evaluation metric & ARI, NMI & Purity, accuracy \\
Training duration & Number of generator calls & Epochs over fixed dataset \\
\bottomrule
\end{tabular}
\end{center}

% ─────────────────────────────────────────────
\section{Non-Trivial Challenges and Limitations of the Naive Approach}
% ─────────────────────────────────────────────

\subsection{The Core Difficulty: Visual Similarity Does Not Equal Source Identity}

The approach outlined in Section~3 is a reasonable starting point, but it rests on an
assumption that deserves scrutiny. Standard contrastive learning methods such as SimCLR
\parencite{simclr2020} work by treating two \emph{augmented views of the same image} as
a positive pair. Because both views are derived from the same image, they look very
similar — and the contrastive loss succeeds by exploiting that visual similarity.

In the present task this assumption breaks down in two directions simultaneously:

\begin{itemize}
    \item \textbf{Intra-image fragments can be visually dissimilar.} A 16$\times$16 tile
          from the sky region of a landscape photograph looks nothing like a tile from the
          ground region of the same photograph. Yet the model is asked to place both close
          together in embedding space.
    \item \textbf{Inter-image fragments can be visually similar.} A sky tile from image A
          and a sky tile from image B may be nearly identical in colour and texture, yet
          the model is asked to push them apart.
\end{itemize}

This is a fundamental tension that a naive pairwise contrastive loss will struggle to
resolve. The model cannot rely on low-level visual similarity as a proxy for source
identity, because that proxy is unreliable in both directions.

\subsection{Implications for Model Design}

This tension suggests that the problem may require approaches beyond standard SimCLR:

\paragraph{Attention mechanisms.}
A transformer-based encoder could allow fragments to attend to each other within a
sample, learning relational features rather than purely local ones. This would give the
model a way to detect global consistency across fragments beyond what any single tile
contains.

\paragraph{Clustering-oriented losses.}
Methods such as SwAV \parencite{swav2020} or SCAN assign soft cluster assignments during
training and use those assignments as part of the loss, making the training objective
more directly aligned with the final clustering goal. These may be better suited to
cases where pairwise visual similarity is not a reliable training signal.

\paragraph{Context as a feature.}
If spatial position is not required as output, it can still be used as an auxiliary
signal during training. Predicting the relative position of two fragments from the same
image is a pretext task \parencite{noroozi2016} that implicitly requires the model to
understand source identity.

\subsection{Summary}

The naive SimCLR approach is a well-justified starting point and may perform adequately
when source images are visually diverse. However, the fundamental mismatch between
visual similarity and source identity means that performance on hard cases — images with
uniform or repetitive content, or datasets with many visually similar images — is likely
to be poor. Recognising this limitation and discussing it explicitly, whether or not it
is fully resolved in the implementation, is an important part of the intellectual
contribution of the work.

\printbibliography

\end{document}
